%% Journal of Supercomputing (Springer) submission
%% Class: svjour3. Download from the Springer author portal and place svjour3.cls here.
%% Compile: pdflatex main -> bibtex main -> pdflatex main x2
%%
%% STATUS: methodology complete, RESULTS ARE PLACEHOLDERS.
%% Every cell marked TBD must be replaced by a number produced by an executed run
%% before this manuscript is submitted anywhere. Do not submit with TBD present.

\documentclass[smallextended]{svjour3}

\usepackage{graphicx}
\usepackage{amsmath}
\usepackage{amssymb}
\usepackage{booktabs}
\usepackage{url}
\usepackage[hidelinks]{hyperref}

\smartqed

\journalname{The Journal of Supercomputing}

\begin{document}

\title{Structure-Exploiting FPGA Acceleration of Superconducting Qubit Hamiltonian
Simulation: A Direct Bosonic Datapath Compared Against Gate-Based Circuit Encoding}

\author{Nasir Ali}

\institute{N. Ali \at
           Embedded Systems Division, Centre for Development of Advanced Computing
           (C-DAC), Noida, India \\
           \email{TBD}}

\date{Received: date / Accepted: date}

\maketitle

\begin{abstract}
Design and calibration of superconducting quantum processors depends on repeated
numerical solution of driven circuit Hamiltonians, and the inner loop of that workload
is latency-bound rather than throughput-bound when it is placed inside a real-time
control path. We present an embedded FPGA architecture for simulating the quantized
LC oscillator and the driven transmon on a Xilinx Alveo U55C, organised around a
structural observation: in their natural bases these Hamiltonians are diagonal plus
bidiagonal (Fock basis) or tridiagonal (charge basis), so a time step costs $O(N)$
rather than the $O(N^2)$ of a general matrix-vector product, and the entire truncated
state fits in on-chip URAM for the truncations of practical interest. We describe a
fully pipelined fixed-point datapath achieving initiation interval one, a build-time
ROM generation flow that ties every coefficient to a physical parameter, and a
quantization-aware machine-learned coarse-step propagator that replaces $K$ fine
Trotter steps with one integer inference. We then implement the same physics a second
way, as a gate-based quantum circuit executed on a native-Verilog FP64 statevector
backend, using both a Pauli-decomposition Trotter construction and a quantum Fourier
transform split-operator construction in which the Josephson nonlinearity reduces to a
single diagonal phase oracle. The two routes are compared at matched accuracy across
wall-clock latency, FPGA resource occupancy, and scaling in the number of coupled
modes, establishing the crossover point at which digitized encoding becomes preferable
to direct propagation.
%% RESULTS SENTENCE: fill in once measured. Example shape, do NOT submit as written:
%% "The direct datapath sustains TBD steps/s at TBD MHz using TBD DSPs, a TBDx
%% reduction in per-step latency relative to a CPU reference at equal fidelity."
\end{abstract}

\keywords{FPGA acceleration \and Quantum simulation \and Superconducting qubits \and
Transmon \and Embedded systems \and Fixed-point arithmetic \and Machine learning
surrogate \and High-level synthesis}

\section{Introduction}
\label{sec:intro}

Superconducting circuits are among the leading platforms for quantum computation, and
their engineering workflow rests on numerical simulation of the underlying circuit
Hamiltonian: extracting transition frequencies and anharmonicities from a proposed
layout, shaping control pulses that suppress leakage out of the computational subspace,
and predicting dispersive readout behaviour \cite{blais2021,krantz2019}. Most of this
work is done offline on CPUs or GPUs, where throughput matters and latency does not.

The situation changes when simulation moves inside a control loop. Real-time
calibration, model-based feedback, and decoder-adjacent estimation all impose
microsecond-scale or tighter deadlines, and at that point the relevant figure of merit
is deterministic per-step latency rather than aggregate FLOPs. This is the regime in
which field-programmable gate arrays are attractive, and it motivates the present work.

Our starting point is the observation that the standard Hamiltonians of circuit quantum
electrodynamics are extremely sparse in the bases in which they are naturally written,
and that this sparsity is structural rather than incidental. Writing the LC oscillator
in second-quantized form gives a purely diagonal Hamiltonian; the capacitively coupled
drive is strictly bidiagonal; the Kerr nonlinearity of the Duffing reduction is again
diagonal; and in the charge basis the Josephson cosine is exactly a shift operator, so
the full transmon Hamiltonian is tridiagonal. A hardware datapath built around these
facts performs $O(N)$ work per time step with all coefficients available from small
on-chip read-only memories, and never needs to materialise a dense operator.

We contribute the following.

\begin{enumerate}
\item A native-Verilog, fully pipelined fixed-point datapath for driven bosonic and
transmon Hamiltonians on an Alveo U55C, with an initiation interval of one and the
truncated state resident in on-chip memory, together with a systematic study of
fixed-point word format against accumulated infidelity.
\item A quantization-aware machine-learned coarse-step propagator, formulated so that
unitarity is enforced structurally rather than by penalty, and evaluated honestly:
we show that for the purely harmonic case the surrogate cannot beat the analytic
propagator and is a methodology testbed only, and that its value appears in the
nonlinear driven transmon regime where no closed form exists.
\item A second, independent implementation of the same physics as a gate-based quantum
circuit, comprising a Pauli-decomposition Trotter construction and a quantum Fourier
transform split-operator construction, executed on an existing native-Verilog FP64
statevector backend.
\item A crossover analysis quantifying, at matched accuracy, the truncation size and
mode count beyond which the digitized circuit route becomes preferable to direct
propagation.
\end{enumerate}

Section~\ref{sec:background} fixes the physics and the numerical formulation.
Section~\ref{sec:direct} describes the direct datapath.
Section~\ref{sec:ml} covers the learned propagator.
Section~\ref{sec:circuit} covers the circuit encodings.
Section~\ref{sec:results} reports measurements and Section~\ref{sec:conclusion}
concludes.

\section{Related work}
\label{sec:related}

%% TODO: cite your own prior work here: the FPGA statevector simulator, the multi-FPGA
%% distributed QFT work, and the HBM2 statevector distribution kernel patent.
%% Placeholder keys are in references.bib and MUST be completed before submission.

Work adjacent to this paper falls into three groups. The first is FPGA acceleration of
quantum circuit simulation, where the dominant concern is the memory-bandwidth cost of
applying gates to an exponentially large state vector, and where high-bandwidth memory
has enabled qubit counts previously restricted to clusters \cite{prior_fpga_sim,
prior_multifpga}. The second is classical simulation of circuit QED itself, typically
in Python frameworks with dense or sparse linear algebra backends \cite{qutip}. The
third is digital quantum simulation of continuous-variable and bosonic systems, whose
theoretical foundations were established in the split-operator constructions of Wiesner
and Zalka and the network formulations of Somma and coworkers \cite{wiesner1996,
zalka1998, somma2002, lloyd1996}.

The present work differs from all three in combining them: we treat the same
Hamiltonian both as an analog-style state-vector propagation problem and as a digitized
circuit, on the same device family, and compare them under a single accuracy metric.

\section{Background and formulation}
\label{sec:background}

\subsection{From the circuit to the oscillator}

For a lumped $LC$ resonator with node flux $\Phi$ and conjugate charge $Q$,
\begin{equation}
H = \frac{Q^2}{2C} + \frac{1}{2} C \omega_{01}^2 \Phi^2,
\qquad
\omega_{01} = \frac{1}{\sqrt{LC}},
\qquad
Z_0 = \sqrt{\frac{L}{C}} .
\label{eq:lc}
\end{equation}
Promoting the conjugate pair to operators and introducing ladder operators through
\begin{equation}
\hat{Q} = i\sqrt{\frac{\hbar}{2 Z_0}}\left(\hat{a}^\dagger - \hat{a}\right),
\qquad
\hat{\Phi} = \sqrt{\frac{\hbar Z_0}{2}}\left(\hat{a}^\dagger + \hat{a}\right),
\label{eq:ladder}
\end{equation}
reduces Eq.~\eqref{eq:lc} to
\begin{equation}
\hat{H} = \hbar \omega_{01}\left(\hat{a}^\dagger \hat{a} + \tfrac{1}{2}\right).
\label{eq:sho}
\end{equation}
The equally spaced spectrum of Eq.~\eqref{eq:sho} is precisely what makes the bare
resonator unusable as a qubit, and motivates the replacement of the linear inductor by
a Josephson junction.

\subsection{The transmon}

With a junction of energy $E_J$ and charging energy $E_C = e^2/2C_\Sigma$,
\begin{equation}
\hat{H} = 4 E_C (\hat{n} - n_g)^2 - E_J \cos\hat{\phi},
\qquad [\hat{\phi}, \hat{n}] = i,
\label{eq:transmon}
\end{equation}
where $n_g$ is the offset charge \cite{koch2007}. In the charge basis truncated to
$|n| \le n_{\max}$ the first term is diagonal and the cosine is exactly the shift
operator,
\begin{equation}
\cos\hat{\phi} = \tfrac{1}{2}\sum_n \left( |n\rangle\langle n{+}1| +
|n{+}1\rangle\langle n| \right),
\label{eq:cosphi}
\end{equation}
so Eq.~\eqref{eq:transmon} is tridiagonal with constant off-diagonal $-E_J/2$. Expanding
the cosine to quartic order gives the Duffing reduction
\begin{equation}
\hat{H} \simeq \hbar\omega_{01}\hat{a}^\dagger\hat{a}
+ \frac{\alpha}{2}\hat{a}^\dagger\hat{a}^\dagger\hat{a}\hat{a},
\qquad
\hbar\omega_{01} = \sqrt{8E_JE_C} - E_C,
\qquad
\hbar\alpha = -E_C .
\label{eq:duffing}
\end{equation}

\subsection{Drive}

The voltage source of the original circuit, coupled through a gate capacitance $C_d$,
contributes $\hat{H}_d(t) = -(C_d/C_\Sigma) V(t) \hat{Q}$ with $V(t) = V_0\cos\omega_d t$.
In the rotating frame and under the rotating-wave approximation this becomes
\begin{equation}
\hat{H} = \hbar\delta\,\hat{a}^\dagger\hat{a}
+ \frac{\hbar\Omega}{2}\left(\hat{a} + \hat{a}^\dagger\right),
\qquad \delta = \omega_{01} - \omega_d .
\label{eq:drive}
\end{equation}
Both the laboratory-frame and rotating-frame forms are implemented; the frame is a
configuration option, not an assumption baked into the datapath.

\subsection{Sparsity, and why it determines the architecture}

In the Fock basis $\hat{a}|n\rangle = \sqrt{n}\,|n-1\rangle$ and
$\hat{a}^\dagger|n\rangle = \sqrt{n+1}\,|n+1\rangle$. Consequently:
Eq.~\eqref{eq:sho} is diagonal, Eq.~\eqref{eq:drive} is bidiagonal, the Kerr term of
Eq.~\eqref{eq:duffing} is diagonal, and Eq.~\eqref{eq:transmon} in the charge basis is
tridiagonal. Every operator application is therefore $O(N)$ with coefficients drawn from
tables of $\sqrt{n}$ and $n(n-1)/2$. A further consequence is used throughout as a
correctness oracle: evolution under the number operator alone is an exact diagonal phase
rotation, leaving all Fock populations invariant.

\subsection{Time discretisation}

The baseline integrator is a second-order split-step,
\begin{equation}
U(\Delta t) \simeq
e^{-i\hat{H}_{\text{diag}}\Delta t/2\hbar}\,
e^{-i\hat{H}_{\text{drive}}\Delta t/\hbar}\,
e^{-i\hat{H}_{\text{diag}}\Delta t/2\hbar},
\label{eq:trotter}
\end{equation}
with a fourth-order Runge-Kutta path retained for cross-checking. Trotter error scaling
follows the standard commutator bounds \cite{childs2021}.

\section{Direct FPGA datapath}
\label{sec:direct}

\subsection{Architecture}

%% TODO: replace with the as-built description and a block diagram once M2 closes.
The kernel is written in native Verilog rather than generated through high-level
synthesis, so that pipeline structure and memory partitioning are controlled explicitly.
The truncated state vector is held in on-chip URAM and BRAM; for the truncations of
interest ($N \le 64$) it does not approach HBM capacity, and HBM is reserved for
streaming trajectory snapshots to the host and for feeding batched parameter sweeps.
Coefficient tables are generated at build time from the configuration file, so that no
constant in the bitstream is disconnected from a physical parameter.

\subsection{Fixed-point format}

Amplitudes are stored in a parameterised fixed-point format, swept over Q1.15, Q2.14 and
Q4.12. Table~\ref{tab:fixedpoint} reports accumulated infidelity against the double
precision golden model, from a rounding-accurate Python emulation of the per-step
phase-rotation recurrence (\texttt{bench/quantization/sweep\_precision.py}, M1; log
committed at \texttt{bench/quantization/results\_quantization\_sweep.json}). This
software emulation is not a substitute for the RTL cosim (V5, M2); resource cost and
$f_{\max}$ are synthesis results and remain TBD until M2 runs through Vivado.
Notably, even Q1.15 accumulates $O(10^{-3})$ infidelity by $10^4$ steps because the
ROM-quantized rotator carries a small systematic per-step angular error that compounds
linearly with step count — Q4.12 accumulates enough phase drift to be unusable at this
horizon, which argues against using it in M2 without periodic phase resynchronisation.

\begin{table}[t]
\caption{Fixed-point format against accumulated infidelity and resource cost,
$N = 32$, $10^4$ steps.}
\label{tab:fixedpoint}
\begin{tabular}{lcccc}
\toprule
Format & Infidelity & DSP & LUT & $f_{\max}$ (MHz) \\
\midrule
Q1.15 & $5.39\times10^{-3}$ & TBD & TBD & TBD \\
Q2.14 & $4.17\times10^{-2}$ & TBD & TBD & TBD \\
Q4.12 & $4.98\times10^{-1}$ & TBD & TBD & TBD \\
\bottomrule
\end{tabular}
\end{table}

\subsection{Parallel engines}

$P$ independent oscillator engines are instantiated to serve parameter sweeps over drive
amplitude and detuning, one trajectory per engine. Throughput in trajectories per second
is reported in Section~\ref{sec:results}.

\section{Learned coarse-step propagator}
\label{sec:ml}

\subsection{Formulation}

We train a map $U_\theta$ taking a state and drive parameters to the state after
$\Delta t = K\,\delta t$ with $K \ge 50$, so that a single inference replaces $K$
applications of Eq.~\eqref{eq:trotter}. Two parameterisations are compared: direct
prediction of the $2N$ real components of the propagated state, and prediction of an
anti-Hermitian generator with the bidiagonal sparsity pattern of the true Hamiltonian,
which enforces norm preservation structurally rather than through a penalty term. The
training objective combines state infidelity with norm and energy-drift penalties.

\subsection{Quantization}

Weights are reduced to eight-bit integers and activations to sixteen bits through
quantization-aware training. Fidelity loss attributable to quantization is reported
separately from fidelity loss attributable to the surrogate approximation itself, since
conflating them obscures which of the two limits deployment.

\subsection{Scope of the claim}

For the harmonic Hamiltonian of Eq.~\eqref{eq:sho} the exact propagator is a closed-form
diagonal phase rotation, and no learned surrogate can improve on it. We therefore make
no speedup claim in that regime and use it purely to validate the training and
quantization pipeline against an exactly known answer. The operative claim concerns the
driven transmon and Duffing systems of Eqs.~\eqref{eq:transmon} and~\eqref{eq:duffing},
where no closed form exists, and it is stated at matched fidelity against the fixed-point
Trotter datapath.

\section{Gate-based circuit encoding}
\label{sec:circuit}

\subsection{Encoding}

A $d$-level mode is mapped onto $m$ qubits. We compare a compact binary encoding with
$d = 2^m$, a Gray-code encoding chosen to reduce the Pauli weight of the ladder
operators, and a one-hot encoding used only as a debugging reference.

\subsection{Construction A: Pauli decomposition and Trotter}

The encoded Hamiltonian is decomposed into Pauli strings and compiled through
CNOT-ladder and $R_Z$ blocks. In the binary encoding the number operator is
\begin{equation}
\hat{n} = \sum_{k=0}^{m-1} 2^k \frac{I - Z_k}{2},
\label{eq:numberop}
\end{equation}
so evolution under Eq.~\eqref{eq:sho} is exact, carries no Trotter error, and costs
exactly $m$ single-qubit $R_Z$ rotations. The diagonal Kerr term likewise compiles to
$R_Z$ and $R_{ZZ}$ rotations only. All Trotter error in the harmonic case therefore
originates in the drive term, which provides a sharp diagnostic.

\subsection{Construction B: quantum Fourier transform split-operator}

Charge and phase are conjugate, and on a truncated periodic charge lattice the
transformation between them is the exact quantum Fourier transform rather than an
approximation to it \cite{wiesner1996, zalka1998}. One step of the transmon evolution is
therefore
\begin{equation}
U(\delta t) =
\mathrm{QFT}^{-1}\,
e^{\,i E_J \cos\hat{\phi}\,\delta t/\hbar}\,
\mathrm{QFT}\,
e^{-i 4 E_C \hat{n}^2 \delta t/\hbar},
\label{eq:splitop}
\end{equation}
in which each exponential is diagonal in its own basis and compiles to $R_Z$ and
controlled-phase rotations, and the entire Josephson nonlinearity is expressed as a
single diagonal phase oracle. The quadratic term is expanded pairwise using
Eq.~\eqref{eq:numberop}.

\subsection{Backend}

Circuits are executed on an existing native-Verilog FP64 statevector simulator for the
Alveo U55C, exposed through a Qiskit \texttt{BackendV2} interface, supporting up to
29 qubits across eight HBM channels. Every circuit is compiled to the backend's native
gate basis, with the cost of any required decomposition recorded rather than assumed
away.

\subsection{Multi-mode extension}

Adding a second mode with Jaynes-Cummings coupling
$\hat{H}_{\text{int}} = \hbar g(\hat{a}^\dagger\hat{b} + \hat{a}\hat{b}^\dagger)$
yields a transmon coupled to a readout resonator, from which the dispersive shift $\chi$
is extracted and pointer-state separation observed. This is the regime in which the
encoding overhead of the circuit route is expected to be repaid.

\section{Results}
\label{sec:results}

%% ============================================================================
%% NOTHING IN THIS SECTION HAS BEEN MEASURED YET.
%% Do not populate any cell from an estimate, an extrapolation, or a prior project.
%% Each number must come from a run whose log is committed under bench/.
%% ============================================================================

\subsection{Validation}

Table~\ref{tab:validation} lists the acceptance gates and their status. V1-V4 (the M0
golden model) pass under \texttt{physics/tests/}. V5 was run in reduced scope (free
evolution only; the RWA drive term is not yet in RTL) via Verilator, standing in for
Vivado xsim, which was not available in the environment this manuscript was drafted in
--- CLAUDE.md's own repo layout names Verilator as an acceptable cosim tool alongside
xsim. It currently fails: at $N=32$, $10^4$ steps, measured infidelity is $4.1\times
10^{-2}$ at Q2.14 (the configured default) and $5.4\times10^{-3}$ at Q1.15, both far
above the $10^{-4}$ gate, from a systematic per-step rotator-quantization error that
compounds over the step count. A software-only re-sweep (same recurrence, not yet built
in RTL) suggests roughly 20 fractional bits are needed, beyond any of the three 16-bit
formats this section commits to --- see \texttt{docs/validation.md} for the full
account. Everything else needs RTL/hardware not yet built, or the toolchain/hardware
access this project has not had, and remains TBD rather than estimated.

\begin{table}[t]
\caption{Validation gates.}
\label{tab:validation}
\begin{tabular}{clc}
\toprule
Gate & Criterion & Status \\
\midrule
V1  & Spectrum matches $\hbar\omega_{01}(n+\tfrac12)$ to $10^{-10}$ & PASS \\
V2  & Free evolution preserves Fock populations exactly & PASS \\
V3  & Coherent state follows the classical trajectory, $Q_M = 0$ & PASS \\
V4  & Resonant drive gives $|\alpha(t)| = \Omega t/2$ & PASS \\
V5  & RTL against golden, infidelity $< 10^{-4}$ at $10^4$ steps & FAIL$^\dagger$ \\
V6  & Surrogate infidelity $< 10^{-3}$ on held-out trajectories & TBD \\
V8  & Encoded free evolution exact to FP64 & TBD \\
V9  & Constructions A and B agree with direct propagation & TBD \\
V10 & Trotter error exponent confirmed by fit & TBD \\
V12 & $E_{01}$ within 1\% of $\sqrt{8E_JE_C} - E_C$ & TBD \\
V13 & DRAG suppresses $|2\rangle$ leakage by $\ge 10\times$ & TBD \\
\bottomrule
\end{tabular}
\end{table}

$^\dagger$Reduced scope: free evolution only, at Q2.14; see text above.

\subsection{Resource utilisation}

\begin{table}[t]
\caption{Resource utilisation by configuration on the Alveo U55C.}
\label{tab:resources}
\begin{tabular}{lccccccc}
\toprule
$N$ & Format & $P$ & ML & LUT & FF & DSP & BRAM/URAM \\
\midrule
32 & Q2.14 & 1  & off & TBD & TBD & TBD & TBD \\
32 & Q2.14 & 1  & on  & TBD & TBD & TBD & TBD \\
32 & Q2.14 & 8  & off & TBD & TBD & TBD & TBD \\
64 & Q2.14 & 1  & off & TBD & TBD & TBD & TBD \\
\bottomrule
\end{tabular}
\end{table}

\subsection{Crossover analysis}

\begin{table}[t]
\caption{Direct datapath against circuit encoding at matched infidelity.}
\label{tab:crossover}
\begin{tabular}{lccccc}
\toprule
Route & Qubits & Depth & 2Q gates & Latency/step & Infidelity \\
\midrule
Direct bosonic datapath      & --- & --- & --- & TBD & TBD \\
Construction A (Pauli/Trotter) & TBD & TBD & TBD & TBD & TBD \\
Construction B (QFT split-op)  & TBD & TBD & TBD & TBD & TBD \\
CPU reference (QuTiP)          & --- & --- & --- & TBD & TBD \\
GPU reference                  & --- & --- & --- & TBD & TBD \\
\bottomrule
\end{tabular}
\end{table}

\subsection{Discussion}

%% Write only after the tables above are filled. If the circuit route loses for a
%% single mode, report that plainly and locate the mode count at which it wins.

\section{Conclusion}
\label{sec:conclusion}

%% Write last.

\begin{acknowledgements}
%% TODO: National Quantum Mission / MeitY acknowledgement text, and colleagues.
\end{acknowledgements}

\section*{Declarations}

\textbf{Funding.} TBD.
\textbf{Conflict of interest.} The author declares no competing interests.
\textbf{Data availability.} All source code, testbenches, and benchmark logs are
available at \url{https://github.com/nasir26/qtransmon-fpga}.

\bibliographystyle{spmpsci}
\bibliography{references}

\end{document}
